% ==========================================================================
% Paper 2 — Section 1: Introduction. The Conflation Problem and the
% Architectural Stance.
% All quantitative statements in this section are forward pointers to the
% frozen archives; the numbers themselves are stated and sourced in the
% sections that own them (§3–§8).
% ==========================================================================
\section{Introduction: The Conflation Problem and the Architectural Stance}
\label{sec:intro}

\subsection{The conflation problem}
\label{sec:conflation}

A striking feature of the modern literature on attention, memory, and
relational computation is that its central vocabulary is used
interchangeably across levels that are causally distinct. A network that is
``attentive'' is described as routing information; a unit whose output is a
weighted mixture of stored items is described as binding them; a system that
can retrieve two stored values is described as composing them. Underneath
this vocabulary sits a chain of five primitives---\emph{memory},
\emph{saliency weighting}, \emph{content routing}, \emph{semantic binding},
and \emph{composition}---each of which imposes its own, strictly stronger,
requirement on the substrate that realizes it. When these are conflated, the
classic inferential error follows: observing the \emph{lower} primitive,
together with rich nonlinear dynamics, is taken as evidence for the
\emph{higher} one. Four examples of the conflation recur in our own field of
view: (a) \emph{history dependence $\Rightarrow$ attention routing}---a
non-Markovian hidden state does not imply that a query can reorder candidate
preferences; (b) \emph{attention weighting $\Rightarrow$ semantic
binding}---a softmax mixture over stored items delivers a convex combination,
not a bound symbol; (c) \emph{sharpening $\Rightarrow$ selection}---any finite
temperature keeps every candidate weight strictly positive, so sharpening
approaches, but never achieves, discrete delivery; and (d) \emph{retrieval
$\Rightarrow$ composition}---delivering two values into a pair is not the
same as applying a deterministic algebraic relation to them.

This is not merely a semantic complaint. End-to-end training of monolithic
networks makes the conflation unfalsifiable in practice: every stage of a
black-box system co-varies under optimization, so post hoc attribution of a
capability to a mechanism is, in general, impossible \citep{fodor1988,
lake2017,lake2018}. The binding problem of classical neuroscience is the
canonical instance: synchrony, assembly formation, and indexing have all been
proposed as binding mechanisms \citep{von1995,singer1999}, yet a demonstration
that a candidate mechanism \emph{binds}, as opposed to merely re-weights, is
rarely separated from the demonstration that it \emph{routes}. What the field
lacks is not more powerful models but a \emph{boundary map}: for each rung of
the primitive ladder, the minimal structural condition under which the
primitive becomes observable in a dynamical substrate, and a falsifiable test
that separates it from the rungs below.

\subsection{Thesis and methodological stance}
\label{sec:stance}

The thesis of this paper is architectural rather than dynamical:

\begin{quote}
\emph{Computational capability is constrained not only by nonlinear
dynamics, but by the organizational degrees of freedom available to the
dynamics.}
\end{quote}

More dynamics---richer nonlinearities, larger state spaces, longer memory
windows---does not, by itself, produce new computational primitives. It
provides a richer \emph{stage}; the primitives are provided by the
\emph{topological organization} of the computation. To make the thesis
testable we adopt two disciplines that run through the entire study.

\emph{Radical minimalism.} The substrate is a single dynamical unit---the
Temporal Attention Neuron (\TAN) of the frozen TAN-I archive---a leaky,
spiking neuron whose input current is a surprise-gated, temporally attended
mixture of its recent inputs (§2.3). It is deliberately the smallest system
in which the conflation (a)--(d) can be exhibited, because every confound
that a larger network would introduce (depth, width, training, shared
substrates) is absent by construction.

\emph{The single-degree-of-freedom ladder.} Each step of the study changes
exactly one organizational degree of freedom and leaves the frozen dynamics,
window, gates, and readouts untouched (§2.4). After each step, the new
configuration is audited by counterfactual experiments---identical histories
with only the query changed---and by negative-control chains in which the
preceding rung of the ladder is shown to fail the same test (§2.5). A rung is
claimed only if the audit passes under the joint criteria, never from a
single scalar error.

We state the resulting scope preemptively. This paper does not propose a
network that beats any benchmark; it proposes a \emph{constraint map} for
dynamical computation. Its results are \emph{sufficiency} statements within
one explicitly frozen mechanism family, not necessity statements about
brains or about all possible machines; and negative results are reported
with the same status as positive ones. Readers expecting accuracy tables
will find boundary tables instead; that is the intended product.

\subsection{Why a single neuron is the right test bed}
\label{sec:testbed}

The TAN-I programme established the bottom two rungs of the ladder inside
this substrate, and---decisively---located the first wall above them. In
brief (full numbers and protocols in §3--§5): the scalar TAN is
historically dependent and its membrane potential is not a sufficient
Markov state (Phase 1: $N=160{,}000$ collision experiments; state
insufficiency lemma), so \emph{memory} is realized. Its surprise-gated
attention reshapes the event-locked response geometry of the system
(Phase 2: effective rank $d_D\approx2.089$ under attention versus
$\approx1.484$ without), so \emph{saliency weighting} is realized and is
geometrically effective. Yet two negative probes froze the boundary:
a temporal distractor task showed no robustness premium attributable to
attention, and a formal binding audit---the query-conditioned
identifiability audit, archived as \textsc{terminated /
audit\_failed}---recorded a counterfactual argmax flip rate of exactly
$0.000$: the scalar kernel $E_{t,i}=\beta W_qW_k\,S_t\,x_i$ can rescale
candidate logits but can never reorder them, because its ``query'' is the
non-negative scalar surprise cone $\mathcal{Q}_+=[0,\infty)$. Scalar
saliency, in other words, lies \emph{below} routing on the ladder, and the
gap is provable rather than empirical (the order-conservation theorem,
§5). The TAN therefore leaves us exactly where a mechanism study should
start: at a wall, with a clean statement of which side each frozen result
belongs to.

\subsection{Contributions and roadmap}
\label{sec:contrib}

The present paper (TAN-II, sprints 4.1--4.3-A) ascends the ladder one degree
of freedom at a time and reports the following results.

\begin{enumerate}
\item \textbf{Competition without routing (§4).} Adding
excitatory--inhibitory opponent topology to the frozen neuron produces
nonmonotonic selective tuning by analytic construction---the E-I
nonmonotonicity theorem---and the static opponent control shows the peak
exists with \emph{no} attention at all ($x^*=1.00$, $h_E=0.200$), severing
the attribution of local tuning to attention. Attention is thereby redefined
as dynamic gain and operating-range modulation. Across six layouts the
counterfactual flip rate remains exactly $0.000$:
$\text{Competition}\not\Rightarrow\text{Routing}$.
\item \textbf{The geometric boundary of routing (§5).} Two-dimensional
vector query--key interaction unlocks routing capacity with an exact law,
$\mathrm{FlipRate}_{\mathrm{CF}}(\theta)=\theta/\pi$ (132/132 Monte Carlo
conditions; winner measures $1/2$; the winner-variation statistic $F_d=1/2$
constant across $d\in\{2,3,4,8\}$ as a negative control). In the running
neuron, changing only the query from $3$ to $4$ on a fixed history flips the
attention winner with margins $+0.2587\to-0.1559$ (closed-form crossing
$Q^*=3.7071$; ensemble flip rate $0.4425$ vs.\ quadrature benchmark
$0.4445$), and key normalization is shown to be a mechanistic enabler, not a
cosmetic control.
\item \textbf{From routing to binding, and to exactness (§6--§7).}
Key--value decoupling yields soft binding whose value-transfer is
winner-aligned and provably pairing-sensitive ($T_{\mathrm{swap}}=-T$ to
$10^{-12}$), with a softness distance $D\approx1.7$--$1.8$ inherited from
softmax convexity. The winner-take-all limit $\gamma\to\infty$ achieves
exact, lossless delivery while the routing statistics
($\mathrm{FlipRate}\approx0.4445$) remain exactly flat across the entire
sharpening ladder---the first orthogonal axis:
\emph{routing selectivity $\perp$ readout discreteness}.
\item \textbf{From binding to composition (§8).} A single head is a
bottleneck: it delivers one scalar symbol per instant and cannot simultaneously
deliver two distinct bound operands to general downstream algebra, so binding
alone cannot compose general arithmetic ($\text{Binding}\not\Rightarrow\text{Composition}$;
the three-error joint audit catches the $V_A=0$ accidental-zero trap). Two parallel channels
plus one minimal two-input algebraic node achieve composition with
\emph{zero} error conditioned on correct binding, over $3\times10^4$
Monte Carlo histories ($E[E_{\mathrm{comp}}|\text{both-routed}]\equiv0$,
$\max=0$), while all unconditional error is attributable to upstream
routing---the second orthogonal axis:
\emph{binding fidelity $\perp$ combiner error}.
\end{enumerate}

These results assemble into the unified mechanistic ladder of §9 and the
graded claim ledger of §10. The through-line of the paper is that each rung
corresponds to one \emph{named} organizational degree of freedom---opponent
topology, vector interaction, key--value decoupling, discrete selection,
parallel channels, and an algebraic node---and that each rung can be
assigned a falsifiable boundary test. More dynamics did not have to be
added at any step; organization did.

The remainder of the paper is organized as follows. §2 gives the formal
framework: axiomatic definitions of the five primitives, the frozen
baseline model, the catalog of extension operators, and the audit
methodology. §3 records the state and geometry baseline inherited from
TAN-I. §4–§8 present the five rungs and their audits in ladder order.
§9 assembles the master table, discusses the biological reading of each
degree of freedom (opponent microcircuits, oscillatory phase, dendritic
subunits, axonal threshold integration), and states the general discussion.
§10 closes with the claim ledger and the reproducibility statement.
